\documentclass[12pt,a4paper]{article}
\usepackage[utf8]{inputenc}
\usepackage{geometry}
\geometry{a4paper, margin=1in}
\usepackage{graphicx}
\usepackage{hyperref}
\usepackage{listings}
\usepackage{xcolor}

\title{Project 15: Real-Time Security Log Visualizer}
\author{AIT Cybersecurity Class Project}
\date{\today}

\begin{document}

\maketitle

\section{Problem and Background}
With the increasing number of cyber threats, security operations centers (SOC) and system administrators require real-time visibility into server access logs to detect malicious activities promptly. Without automated parsing and visualization, identifying brute-force attacks or suspicious off-hour logins from raw log files is a tedious and error-prone process. This project addresses the need for a localized, lightweight, real-time threat detection and visualization tool.

\section{Objective}
The primary objective is to build a Real-Time Security Log Visualizer that monitors a local \texttt{access.log} file in real-time. The application parses newly appended log entries, detects specific security threats (such as brute-force attacks and off-hour logins), and updates a dynamic dashboard to provide immediate situational awareness.

\section{Methodology and Approach}
We adopted a File Watcher methodology (Method 2) to observe log file modifications continuously. An asynchronous background process monitors the target file for changes and triggers a parsing routine whenever new data is appended. Data is retained in memory, enabling immediate analysis of historical context, such as calculating the frequency of failed logins per IP address over time. The results are instantly reflected on a web-based dashboard.

\section{Tools and Technologies Used}
\begin{itemize}
    \item \textbf{Python 3}: Core programming language used for scripting the backend logic and data simulation.
    \item \textbf{Streamlit}: A rapid web framework used to build the real-time interactive dashboard, including metrics tables and bar charts.
    \item \textbf{Watchdog}: A Python API and shell utility used to monitor file system events (specifically, modifications to the \texttt{access.log} file).
    \item \textbf{Pandas}: Employed for efficient data manipulation, filtering, and aggregation of the log records.
\end{itemize}

\section{Implementation and Experiment Details}
The implementation consists of two primary components:
\begin{enumerate}
    \item \textbf{Log Simulator (\texttt{logger\_sim.py})}: Generates mock access logs containing timestamps, random IP addresses, and login statuses (Success or Failed). It introduces artificial bias to simulate brute-force attacks from specific IPs and generates off-hour timestamps.
    \item \textbf{Dashboard Application (\texttt{app.py})}: Utilizes Streamlit and Watchdog. It captures new log lines, parses the timestamp and status, and applies the following detection rules:
    \begin{itemize}
        \item \textit{Brute Force}: Flags any IP address appearing more than 5 times with a "Failed Login" status.
        \item \textit{Off-Hours}: Identifies any login attempt occurring outside the standard operational window of 9:00 AM to 6:00 PM.
    \end{itemize}
\end{enumerate}

\section{Results and Findings}
During testing, the Log Simulator correctly streamed data into the target file. The dashboard responded in real-time (sub-second latency) without requiring manual browser refreshes. Brute-force attacks were successfully detected and highlighted as "Critical" threats once the 5-failure threshold was breached. Off-hour login attempts were accurately tagged. The visual metrics provided a clear, real-time overview of the system's security posture.

\section{Limitations}
\begin{itemize}
    \item \textbf{Scalability}: Keeping all logs in memory via Streamlit's session state is sufficient for demonstration but unsuitable for production environments with massive log volumes.
    \item \textbf{Persistence}: The application resets its detected state upon a full server restart because data is not backed by a persistent database.
    \item \textbf{Log Format Rigidness}: The parsing logic relies on a strict formatting structure; any deviation in the log source could break the parser.
\end{itemize}

\section{Lessons Learned}
\begin{itemize}
    \item Integrating asynchronous event-driven libraries (Watchdog) with synchronous web frameworks (Streamlit) requires careful state management to avoid race conditions.
    \item Real-time visual feedback is highly effective for security monitoring compared to static log analysis.
    \item Simulating realistic threat actor behavior is essential for thoroughly testing detection algorithms.
\end{itemize}

\section{Recommendations and Future Improvement}
\begin{itemize}
    \item \textbf{Database Integration}: Transition from in-memory data structures to a robust time-series database (e.g., InfluxDB) or a SIEM backend (e.g., Elasticsearch).
    \item \textbf{Advanced Analytics}: Incorporate machine learning models to identify anomalous login patterns beyond simple threshold-based heuristic rules.
    \item \textbf{Alerting Mechanism}: Add webhook support or email integration to notify administrators automatically when critical thresholds (like brute-force) are met.
\end{itemize}

\end{document}
